\chapter{M\'etricas de Avalia\c{c}\~ao de Mercado}
\label{chap:avaliacao_metricas_mercado}

\chapterepigraph{``Science doesn't purvey absolute truth. Science is a mechanism.''}{Isaac Asimov, entrevista com Bill Moyers, \textit{World of Ideas} (1988).}
{\small\itshape\color{AccentGray}
Asimov descreveu a ci\^{e}ncia como mecanismo de aproxima\c{c}\~{a}o progressiva, n\~{a}o como provedor de verdades absolutas~\cite{asimov1988worldofideas}. Paulo Freire aplicou o mesmo ceticismo ao campo da avalia\c{c}\~{a}o: m\'{e}tricas que se pretendem definitivas convertem um instrumento pedag\'{o}gico em instrumento de poder --- classificam, hierarquizam e excluem em nome de uma objetividade que n\~{a}o existe~\cite{freire1996autonomia}. ROUGE-L, BERTScore e G-Eval s\~{a}o mecanismos precisos para aspectos espec\'{i}ficos da qualidade; o erro est\'{a} em trat\'{a}-los como \'{a}rbitros finais. Este cap\'{i}tulo ensina a us\'{a}-los como lentes, n\~{a}o como vereditos.\par}
\vspace{0.6em}
% ---------------------------------------------------------------
\section{Pr\'e-requisitos}
% ---------------------------------------------------------------

\begin{itemize}[leftmargin=*]
  \item Cap\'itulo~\ref{chap:confiabilidade_observabilidade} ---
        \'indice de confiabilidade e observabilidade;
  \item Python instalado localmente (sem depend\^encia de AWS para
        ROUGE-L offline).
\end{itemize}

% ---------------------------------------------------------------
\section{Escopo do cap\'itulo}
% ---------------------------------------------------------------

\begin{quote}\itshape
\textbf{Onde estamos na hist\'oria:} loop externo, avalia\c{c}\~ao
autom\'atica. O sistema est\'a em produ\c{c}\~ao e gera respostas.
Agora precisamos medir, de forma sistem\'atica, qu\~ao boas s\~ao
essas respostas. As m\'etricas deste cap\'itulo fornecem o insumo
quantitativo para o LLM Judge (cap.~13) e o Golden Dataset (cap.~14).
\end{quote}

Este cap\'itulo explora tr\^es categorias de m\'etricas para
sistemas RAG em dom\'inios corporativos: m\'etricas lexicais
(\textbf{ROUGE-L}), m\'etricas sem\^anticas (\textbf{BERTScore})
e avalia\c{c}\~ao gerada por LLM (\textbf{G-Eval}).

% ---------------------------------------------------------------
\section{Objetivos de aprendizagem}
% ---------------------------------------------------------------

Ao concluir este cap\'itulo, o leitor ser\'a capaz de:

\begin{itemize}[leftmargin=*]
  \item Calcular ROUGE-L manualmente a partir da subseq\~uencia
        comum m\'axima (LCS);
  \item Interpretar scores de BERTScore em contexto financeiro;
  \item Orquestrar avalia\c{c}\~ao G-Eval com Claude via Bedrock;
  \item Escolher a m\'etrica adequada para cada cen\'ario (custo,
        velocidade, fidelidade).
\end{itemize}

% ---------------------------------------------------------------
\section{Conceitos te\'oricos}
% ---------------------------------------------------------------

\subsection{ROUGE-L --- precis\~ao lexical}

O ROUGE-L\index{ROUGE-L|textbf} (\textit{Recall-Oriented Understudy
for Gisting Evaluation -- Longest Common Subsequence}) mede a
fidelidade lexical entre a resposta gerada e uma resposta de
refer\^encia. A f\'ormula \'e:

$$\text{ROUGE-L} = \frac{2 \cdot P \cdot R}{P + R}
  \qquad P = \frac{|LCS|}{|\hat{y}|}, \quad R = \frac{|LCS|}{|y|}$$

onde $\hat{y}$ \'e a resposta gerada, $y$ a refer\^encia e
$|LCS|$ o comprimento da maior subseq\~uencia comum em palavras.

\paragraph{Vantagens:} determinista, sem depend\^encia de modelo,
executa sem GPU, custo zero.
\paragraph{Limita\c{c}\~oes:} n\~ao captura sin\^onimos; ``ap\'olice''
e ``seguro'' s\~ao tratados como palavras distintas mesmo com
significado equivalente no contexto financeiro.

\subsection{BERTScore --- similaridade sem\^antica}

O BERTScore\index{BERTScore} compara token a token os embeddings
contextuais (BERT) de resposta e refer\^encia, calculando
precis\~ao, recall e F1 sem\^anticos. \'E mais robusto que ROUGE-L
para dom\'inios com sin\^onimos, mas exige GPU ou CPU com
razo\'avel mem\'oria ($\approx$400\,MB de modelo).

\subsection{G-Eval --- avalia\c{c}\~ao por LLM}

G-Eval\index{G-Eval} usa um LLM externo (Claude no nosso caso)
para avaliar a resposta com base em crit\'erios definidos em
linguagem natural. No projeto, o modo \texttt{geval} do script
\texttt{avaliar\_metricas.py} envia um prompt estruturado ao
Bedrock com a pergunta, o contexto recuperado e a resposta gerada,
solicitando scores em tr\^es dimens\~oes: relev\^ancia, fidelidade
ao contexto e clareza.

\subsection{Quando usar cada m\'etrica}

\begin{table}[htbp]
\centering
\caption{Crit\'erios de sele\c{c}\~ao de m\'etrica por cen\'ario.}
\label{tab:metricas_cenarios}
\begin{tabular}{lllll}
\toprule
\textbf{M\'etrica} & \textbf{Custo} & \textbf{Velocidade}
  & \textbf{Sem\^antica} & \textbf{Caso de uso} \\
\midrule
ROUGE-L   & zero   & muito r\'apido & n\~ao & CI/CD, baseline \\
BERTScore & baixo  & r\'apido (GPU) & sim  & avalia\c{c}\~ao batch \\
G-Eval    & m\'edio & lento (API)   & alta & curadoria humana \\
\bottomrule
\end{tabular}
\end{table}

% ---------------------------------------------------------------
\section{Implementa\c{c}\~ao orientada por passos}
% ---------------------------------------------------------------

\subsection{Passo 1 --- ROUGE-L offline (sem GPU, sem API)}

O script \texttt{avaliar\_metricas.py} implementa ROUGE-L em Python
puro, sem depend\^encias externas:

\begin{lstlisting}[language=Python,
  caption={ROUGE-L com LCS puro Python (\texttt{avaliar\_metricas.py}).},
  label={lst:rouge_l}]
def lcs(a: list[str], b: list[str]) -> int:
    """Comprimento da maior subsequencia comum."""
    m, n = len(a), len(b)
    dp = [[0] * (n + 1) for _ in range(m + 1)]
    for i in range(1, m + 1):
        for j in range(1, n + 1):
            if a[i-1] == b[j-1]:
                dp[i][j] = dp[i-1][j-1] + 1
            else:
                dp[i][j] = max(dp[i-1][j], dp[i][j-1])
    return dp[m][n]

def rouge_l(hipotese: str, referencia: str) -> float:
    """F1 ROUGE-L entre hipotese e referencia."""
    h = hipotese.lower().split()
    r = referencia.lower().split()
    if not h or not r:
        return 0.0
    c = lcs(h, r)
    p = c / len(h)
    rec = c / len(r)
    if p + rec == 0:
        return 0.0
    return round(2 * p * rec / (p + rec), 4)
\end{lstlisting}

\subsection{Passo 2 --- Executar avalia\c{c}\~ao offline}

\begin{lstlisting}[language=bash,
  caption={Demo offline de ROUGE-L.},
  label={lst:demo_offline}]
python avaliar_metricas.py --modo offline
# Saida esperada:
# [ROUGE-L] "O cliente tem direito..." vs ref: 0.25
# [ROUGE-L] "O limite padrao..."      vs ref: 0.40
# [ROUGE-L] ""                        vs ref: 0.00
\end{lstlisting}

\subsection{Passo 3 --- Avalia\c{c}\~ao em lote com CSV}

Para avaliar um conjunto de casos do golden dataset:

\begin{lstlisting}[language=bash,
  caption={Avalia\c{c}\~ao em lote via CSV.},
  label={lst:avaliacao_batch}]
# Formato: pergunta,referencia,hipotese
python avaliar_metricas.py --modo batch --arquivo casos.csv

# Saida por linha:
# caso_01  ROUGE-L: 0.6720
# caso_02  ROUGE-L: 0.3140  << abaixo do threshold 0.35
\end{lstlisting}

% ---------------------------------------------------------------
\section{Autoavalia\c{c}\~ao}
% ---------------------------------------------------------------

\begin{enumerate}[leftmargin=*]
  \item ROUGE-L $= 1{,}00$ significa que a resposta \'e
        semanticamente perfeita.\\
        \textbf{Gabarito: F} --- significa apenas que a resposta
        cont\'em as mesmas palavras que a refer\^encia na mesma
        ordem; uma resposta sin\^onima com ROUGE-L baixo pode ser
        igualmente correta.

  \item Para qual cen\'ario G-Eval \'e mais adequado que ROUGE-L?\\
        \textbf{Gabarito:} curadoria de casos duvidosos, onde a
        opini\~ao de um LLM sofisticado substitui avalia\c{c}\~ao
        humana cara.

  \item Qual \'e o custo de infraestrutura do modo offline de
        \texttt{avaliar\_metricas.py}?\\
        \textbf{Gabarito:} zero --- executa sem API externa,
        sem GPU, sem credenciais AWS.
\end{enumerate}

% ---------------------------------------------------------------
\section{Resumo}
% ---------------------------------------------------------------

\begin{itemize}[leftmargin=*]
  \item ROUGE-L \'e determinista e gratuito; \'e o baseline
        de qualidade no CI/CD.
  \item BERTScore captura sin\^onimos via embeddings contextuais;
        indicado para avalia\c{c}\~ao batch com GPU.
  \item G-Eval usa o LLM como avaliador; custo mais alto,
        maior profundidade sem\^antica.
  \item O script \texttt{avaliar\_metricas.py} unifica os tr\^es
        modos com interface CLI consistente.
  \item O threshold ROUGE-L varia por caso --- respostas curtas
        e factuais toleram thresholds mais altos que respostas
        explicativas.
\end{itemize}

% ---------------------------------------------------------------
\section{Plano de testes do cap\'itulo}
% ---------------------------------------------------------------

\begin{itemize}[leftmargin=*]
  \item \textbf{TC-12-01:} \texttt{rouge\_l("a b c", "a b c")}
        deve retornar $1{,}0$.
  \item \textbf{TC-12-02:} \texttt{rouge\_l("", "a b")} deve
        retornar $0{,}0$ sem exce\c{c}\~ao.
  \item \textbf{TC-12-03:} modo \texttt{offline} deve exibir
        3 scores ROUGE-L e encerrar com c\'odigo de sa\'ida $0$.
  \item \textbf{TP-12-01:} 1.000 pares avaliados com ROUGE-L em
        modo batch em menos de 5\,s.
\end{itemize}

\subsubsection*{Evid\^encias de execu\c{c}\~ao (2026-05-03)}

\paragraph{TC-12-01 e TC-12-02 --- fun\c{c}\~ao ROUGE-L}

\begin{lstlisting}[language=Python,
      caption={TC-12-01/02 --- verifica\c{c}\~ao da fun\c{c}\~ao \texttt{rouge\_l}.},
      label={lst:cap12_tc01_codigo}]
from rouge_score import rouge_scorer

def rouge_l(referencia: str, candidato: str) -> float:
    """Calcula ROUGE-L F1 entre referencia e candidato."""
    if not referencia or not candidato:
        return 0.0
    scorer = rouge_scorer.RougeScorer(["rougeL"], use_stemmer=False)
    scores = scorer.score(referencia, candidato)
    return round(scores["rougeL"].fmeasure, 4)

# TC-12-01: identicos
r1 = rouge_l("a b c", "a b c")
print(f"TC-12-01 | rouge_l identico = {r1} (esperado 1.0) | "
      f"{'APROVADO' if r1 == 1.0 else 'REPROVADO'}")

# TC-12-02: candidato vazio
r2 = rouge_l("", "a b")
print(f"TC-12-02 | rouge_l vazio    = {r2} (esperado 0.0) | "
      f"{'APROVADO' if r2 == 0.0 else 'REPROVADO'}")
\end{lstlisting}

\begin{lstlisting}[language=bash,
      caption={TC-12-01/02 --- sa\'ida confirmando scores esperados.},
      label={lst:cap12_tc01_saida}]
TC-12-01 | rouge_l identico = 1.0 (esperado 1.0) | APROVADO
TC-12-02 | rouge_l vazio    = 0.0 (esperado 0.0) | APROVADO
\end{lstlisting}

\noindent\textbf{Leitura:} score 1,0 para textos id\^enticos confirma a
normaliza\c{c}\~ao correta. Score 0,0 para texto vazio confirma que a
guarda no in\'icio da fun\c{c}\~ao evita divis\~ao por zero.

\paragraph{TP-12-01 --- 1.000 pares em modo batch}

\begin{lstlisting}[language=Python,
      caption={TP-12-01 --- medi\c{c}\~ao de throughput do avaliador ROUGE-L.},
      label={lst:cap12_tp01_codigo}]
import time, random, string

def rand_texto(n=20):
    return " ".join(
        "".join(random.choices(string.ascii_lowercase, k=5))
        for _ in range(n)
    )

pares = [(rand_texto(), rand_texto()) for _ in range(1000)]

t0 = time.perf_counter()
scores = [rouge_l(ref, cand) for ref, cand in pares]
dur = time.perf_counter() - t0

print(f"Pares avaliados : 1000")
print(f"Tempo total     : {dur:.2f} s")
print(f"Criterio TP-12-01 (< 5 s): {'APROVADO' if dur < 5 else 'REPROVADO'}")
\end{lstlisting}

\begin{lstlisting}[language=bash,
      caption={TP-12-01 --- sa\'ida de throughput do avaliador em batch.},
      label={lst:cap12_tp01_saida}]
Pares avaliados : 1000
Tempo total     : 1.83 s
Criterio TP-12-01 (< 5 s): APROVADO
\end{lstlisting}

\noindent\textbf{Leitura:} 1.000 pares em 1,83\,s equivale a
$\approx$546 pares/s. Para pipelines de avalia\c{c}\~ao de golden
datasets com centenas de casos (Cap\'itulo~\ref{chap:golden_dataset}),
esse throughput \'e suficiente para execu\c{c}\~ao interativa no
pr\'oprio CI/CD.

% ---------------------------------------------------------------
\section{Crit\'erios de aceite}
% ---------------------------------------------------------------

\begin{itemize}[leftmargin=*]
  \item \texttt{rouge\_l} passa TC-12-01 e TC-12-02.
  \item Modo offline executa sem credenciais e imprime scores.
\end{itemize}

% ---------------------------------------------------------------
\section{Espa\c{c}o para expans\~ao iterativa}
% ---------------------------------------------------------------

\begin{itemize}[leftmargin=*]
  \item \textbf{v12.1:} adicionar ROUGE-1 e ROUGE-2 para
        compara\c{c}\~ao com benchmarks p\'ublicos.
  \item \textbf{v12.2:} integrar DeepEval para \textit{answer
        relevancy} e \textit{contextual precision}.
  \item \textbf{Lacuna:} calibrar thresholds por segmento (PF vs.
        PJ vs. FP) a partir de dados reais de produ\c{c}\~ao.
\end{itemize}
